Online marketplaces are central to how communities exchange goods, yet they concentrate trust, custody, and dispute resolution in a single intermediary that charges fees, can fail, and must be relied upon not to misbehave. For a campus community trading low-value items, that intermediary is both costly and unnecessary. We ask whether the marketplace's core trust functions---value, listings, escrow, and reputation---can instead be enforced entirely by public smart contracts, making custody, settlement, and reputation auditable while removing the single point of failure~\cite{nakamoto2008bitcoin,buterin2014ethereum}. This question matters because escrowed peer-to-peer commerce is a broad, recurring pattern (resale, freelance milestones, rental deposits), and a disintermediated, standards-based realization is both practically useful and a clear fit for a public blockchain (Section~\ref{sec:why}).

We present Enigma, a minimal but complete marketplace in which value, listings, escrow, and reputation are all enforced on-chain. Enigma is engineered as four \emph{vertical slices}, each owned by one team member and spanning a contract function set, a front-end module, and a test suite: (1) Token + Wallet, (2) Listings, (3) Escrow and Ratings, and (4) Reputation. This decomposition lets members work and be evaluated independently while composing into one deployed system. Section~\ref{sec:related} reviews related work; Section~\ref{sec:motiv} gives a motivating example and threat model; Section~\ref{sec:why} argues the blockchain fit; Section~\ref{sec:arch} presents the architecture and mitigations; Sections~\ref{sec:slice1}--\ref{sec:slice4} detail the slices; Section~\ref{sec:eval} states the hypothesis, metrics, methodology, and empirical evidence; and Section~\ref{sec:concl} concludes.
